\section{Conclusion}
\label{sec:conclusion}

DenseSLM4MoE demonstrates a practical hybrid architecture achieving GPT-2-scale parameters (1.45B) with significantly improved active-parameter efficiency ($\sim$768M active vs 1.45B total). The combination of MLA (compressed KV), Mamba2 (efficient SSM layers), and aux-loss-free MoE provides a strong foundation for agent workloads at small scale.

\subsection{Current Status}
\begin{itemize}
    \item DenseSLM4MoE v2 pretraining is complete (190{,}967 steps, one epoch over $\sim$6.3B tokens), reaching an evaluation perplexity of 11.49
    \item Architecture is validated through ablation and v1 experiments
    \item Training is stable with no loss spikes or divergence
\end{itemize}

\subsection{What's Needed to Reach ``Useful Agent Level''}

Concrete next steps are detailed in Section~\ref{sec:futurework}.

\subsection{Limitations}
\begin{itemize}
    \item Current benchmarks (C-Eval 28.1\%, TruthfulQA 20\%) are below the threshold for truly useful agents
    \item No RLHF or DPO has been applied yet --- pure next-token prediction pretraining
    \item GSM8K near-random suggests mathematical reasoning needs targeted improvement
    \item Chinese performance bottlenecked by data quality, not architecture
\end{itemize}

\subsection*{Acknowledgments}

We thank the \texttt{fla} \citep{fla} and \texttt{mamba\_ssm} teams for their excellent open-source implementations of Mamba2 and MLA. We also thank the HuggingFace team for LightEval \citep{lighteval}.

\subsection*{Reproducibility}

All model configurations, training scripts, and evaluation pipelines are available at:
\begin{verbatim}
https://github.com/ZiceWang/denseslm4
\end{verbatim}

Training can be launched via:
\begin{verbatim}
python -m denseslm4.train_moe --output_dir runs/denseslm4_moe_v2 ...
\end{verbatim}

Evaluation:
\begin{verbatim}
python -m lighteval --model_args "pretrained=./runs/denseslm4_moe/final_model,..." ...
\end{verbatim}
